\section{G\"odel coding}
\label{sec:arithmetic-coding-04}

\begin{definition}[G\"odel numbering]
A \emph{G\"odel numbering} is an injective effective encoding of the
specified syntactic objects by natural numbers.  Primitive recursiveness of
recognition and of the operations on these codes is a property to be proved,
not an additional hypothesis used to obtain arithmetization.
\end{definition}

We use the following numbering for the primitive arithmetic language
\(\Lang_A=\{0,S,+,\times\}\), with equality logical, and the actual
\(\ND_{\FOL}\) calculus of \cref{sec:fol-deduction-01} and
\cref{sheet:nd-fol}.  Thus \(T\derives\sigma\) in this section means
\(T\derives_{\ND_{\FOL}}\sigma\).  This fixes the effective-calculus
choice in the earlier use of \(\derives\); it includes the classical rule
\(\mathsf{RAA}\), equality, and explicit alpha conversion.  We encode
fully expanded formulas: numerals, infix function notation, comparisons,
and marked bounded quantifiers are the previously declared abbreviations.
The program numbering and computation-history coding of Part IV are unchanged.
Write \(\ulcorner E\urcorner\) for the code in its declared syntactic
sort; \(\overline{\ulcorner E\urcorner}\) is its arithmetic numeral.

\paragraph{Lists and constructor codes.}
Put
\[
 C(a,b)=1+\operatorname{pair}_{\N}(a,b),\qquad
 []=0,\qquad [a_0,\ldots,a_{k-1}]=C(a_0,[a_1,\ldots,a_{k-1}]).
\]
Here \(\operatorname{pair}_{\N}(a,b)=(a+b)(a+b+1)/2+b\) is the
already fixed numerical pairing.  For \(c>0\), let \(h(c),t(c)\)
be the unique pair with \(C(h(c),t(c))=c\), and put
\(h(0)=t(0)=0\).  Both components are less than \(c\).
Consequently every natural number is the code of exactly one finite list;
zero codes its empty list.  Define its length and entries by
\[
 \ell(c)=\min\{j\le c:t^j(c)=0\},\qquad
 (c)_j=\begin{cases}h(t^j(c)),&j<\ell(c),\\0,&j\ge\ell(c).
 \end{cases}
\]
The superscript here means iteration.  Define
\(N_k(a_0,\ldots,a_{r-1})=C(k,[a_0,\ldots,a_{r-1}])\).
The outer constructor tag and its argument list are uniquely recoverable.

The following table is the complete tree grammar; each child displayed as a
term or formula must itself have that type, and the argument list must have
exactly the displayed length.  Variable indices \(i\) range over \(\N\).
\[
\begin{array}{c|c@{\qquad}c|c}
 \text{object}&\text{code}&\text{object}&\text{code}\\\hline
 x_i&N_0(i)&0&N_1()\\
 S(r)&N_2(\ulcorner r\urcorner)&+(r,t)&N_3(\ulcorner r\urcorner,\ulcorner t\urcorner)\\
 \times(r,t)&N_4(\ulcorner r\urcorner,\ulcorner t\urcorner)&r=t&N_5(\ulcorner r\urcorner,\ulcorner t\urcorner)\\
 \bottom&N_6()&\neg\varphi&N_7(\ulcorner\varphi\urcorner)\\
 \varphi\land\psi&N_8(\ulcorner\varphi\urcorner,\ulcorner\psi\urcorner)&
 \varphi\lor\psi&N_9(\ulcorner\varphi\urcorner,\ulcorner\psi\urcorner)\\
 \varphi\to\psi&N_{10}(\ulcorner\varphi\urcorner,\ulcorner\psi\urcorner)&
 \forall x_i\varphi&N_{11}(i,\ulcorner\varphi\urcorner)\\
 \exists x_i\varphi&N_{12}(i,\ulcorner\varphi\urcorner)&&
\end{array}
\]
Each proper child code is less than its parent.  Zero, a wrong tag, a wrong
arity, or a child of the wrong type is not a term or formula code.  In
particular a nullary function and the empty list are different objects.

For literal symbols assign token numbers \(0,\ldots,14\), in this order,
to \(0,S,+,\times,=,\bottom,\neg,\land,\lor,\to,\forall,\exists,(,),\text{comma}\);
assign \(15+i\) to \(x_i\).  A symbol has code \(N_{13}(a)\), where
\(a\) is its token number; a string of tokens \(a_0\cdots a_{k-1}\)
has code \(N_{14}([a_0,\ldots,a_{k-1}])\).  These sorts are disjoint
from tree codes.  A formula's governing code is its tree code, not its
string code.  Its canonical spelling is exactly the prefix term and
parenthesized formula grammar of \cref{ch:fol-syntax-semantics}:
\(S(r),+(r,t),\times(r,t),r=t,\bottom,\neg\varphi,
(\varphi\mathbin\circ\psi),Qx_i\varphi\).
Spaces are not tokens.  The empty string is a valid string but no expression.
For a symbol code require outer tag \(13\) and one argument, then use the
finite token table to test a logical, punctuation, or function-symbol kind;
a token at least \(15\) is a variable, with index its token minus \(15\).
For a string require outer tag \(14\) and one list argument.  Constructor
injectivity and induction on tree formation prove injectivity in each
syntactic sort; the distinct outer tags separate symbols and strings from
terms and formulas.

\paragraph{Primitive-recursive implementation of the coding.}
The following elementary closure details specify the bounds used below.
Truncated subtraction is obtained by recursion using predecessor; zero tests,
order tests, and equality tests follow from it.  Bounded sums and products
are primitive recursions, and a bounded least search is a bounded sum using
the tests that all earlier candidates fail (with a specified default on
failure).  These observations use precisely the class defined in
\cref{sec:turing-machines-02}.  The triangular numbers satisfy
\(u(0)=0, u(n+1)=u(n)+n+1\), so the displayed pairing is primitive
recursive without adding division as an initial operation.  Its inverse
uses two searches bounded by \(c\) for \(C(a,b)=c\).  Iteration of
\(t\), length, entry, and all fixed-length constructor projections are
therefore primitive recursive.

Reversing a list repeatedly removes its head and prepends it to an
accumulator, for exactly \(\ell(c)\) iterations; concatenation reverses
the first list and prepends its entries to the second.  Filtering, bounded
membership, and finite union have the same bounded-loop form.  A loop whose
state is a coded list, whose single update is primitive recursive, and
whose iteration bound is a primitive-recursive function of its inputs is
primitive recursive: first iterate its update by primitive recursion on a
separate counter, then substitute the bound.  No bound on the numerical
size of intermediate states is needed for this closure argument.

Course-of-values computations on codes are implemented by a list of the
already computed rows \(0,\ldots,c-1\).  At step \(c\), earlier rows are
retrieved by bounded list access and the new row is appended.  Thus any
specified primitive-recursive update using only proper child codes defines
a primitive-recursive function.  Apply this simultaneously to the two type
tests in the table: at row \(c\) test its tag and arity and consult the
required earlier term/formula tests.  This proves recognition, rather than
assuming it.  The same construction computes canonical spelling, tree-node
count, and lists of all and of free variable indices.  At \(Qx_i\varphi\),
the all-variable list includes \(i\), and the free-variable list removes
\(i\) from that of \(\varphi\).  A sentence is a formula with empty
free-variable list.  All these functions return zero on an invalid input
tree; their uses always include the appropriate type test.

For completeness, recognition starting from a string also has an explicit
primitive-recursive bound.  Set
\[
 L_0(b)=0,\quad L_{j+1}(b)=C(b,L_j(b)),\quad
 B(0,q)=q+16,\quad
 B(k+1,q)=B(k,q)+1+C(14,L_2(B(k,q))).
\]
If a string has \(k\) tokens, all at most \(q\), its expression tree,
if any, has at most \(k\) nodes and code at most \(B(k,q)\).
Indeed variable indices are at most \(q\), tags are at most \(12\),
arities at most two, and an argument list of length at most two with
entries at most \(b\) is at most \(L_2(b)\).  Induction on tree height,
using monotonicity of pairing, gives the bound.  Search all tree codes up
to \(B(k,q)\) for the required type and identical canonical spelling.
Unique readability makes the result unique.  Use \(q=0\) for the empty
string and return zero when the search fails.  This supplies primitive-
recursive conversion in both directions, with explicit malformed cases.

\paragraph{Substitution and alpha conversion.}
Let \(\nu(0)=N_1()\) and \(\nu(n+1)=N_2(\nu(n))\); thus
\(\nu(n)=\ulcorner\overline n\urcorner\) by induction on \(n\).
For numeral substitution in a tree, replace \(N_0(i)\) by \(\nu(n)\),
leave other variables and \(0,\bottom\) unchanged, and commute with every
function, equality, and Boolean constructor.  At a quantifier with binder
\(i\), leave its entire subtree unchanged; at any other binder recur on
the matrix without renaming.  A numeral has no variables, so these are
exactly all applicable clauses of the book's capture-avoiding substitution.
The course-of-values construction just given proves this total function
primitive recursive, with value zero on malformed input.

The rule checker also needs arbitrary term substitution.  Reject an input
that is not a term or formula, or a replacement that is not a term, and
return zero in those cases.  Use the precise
three quantifier clauses of \cref{sec:fol-syntax-semantics-04}, including
the least fresh index when renaming is required.  Here is a bounded
implementation, which also proves primitive recursiveness of that operation.
Traverse the input using a stack of constructor frames, processing children
from left to right and assembling the result when their frames return.
An entry frame records the current tree, substitution index, replacement
term, and mode (main substitution or fresh renaming).  A return frame
records the constructor, its binder index if any, its unprocessed children,
and the already returned child codes.  Both are fixed-field list records;
the stack is a list of them.  Entering a nonterminal constructor replaces
its entry by a return frame and the first child's entry.  Returning a
child appends its result to that return frame and either enters the next
child or returns the constructor assembled from the completed list.  A
terminal return propagates its code to the waiting frame; an empty waiting
stack marks completion.  These operations specify the traversal update
using only the list operations already constructed.
At a variable return the replacement term or the unchanged variable; the
constant and constructor cases are as above.  At \(Qx_j\psi\), stop
if \(j=i\).  If \(j\ne i\) and \(j\) is absent from the replacement
term, process \(\psi\).  Otherwise find the least index \(z\) absent
from \(\operatorname{Var}(\psi)\cup\operatorname{Var}(t)\cup\{i,j\}\),
first replace the free occurrences of \(x_j\) in \(\psi\) by \(x_z\),
and then process that renamed matrix with the original replacement request.
For this case a continuation frame saves the original request and binder
\(z\), enters the preliminary renaming mode, and on its return enters
main mode on the returned matrix before assembling that quantifier.
The search for \(z\) is bounded by one plus the largest index in that
finite union.  The preliminary renaming stops beneath another \(x_j\)
binder.  Because \(z\) is absent from the whole matrix, this preliminary
pass requires no further renaming, and changes no tree shape.

Every main-traversal call concerns one node of the original tree with only
variable names possibly changed; an inserted copy of \(t\) is returned
as a code and is not processed again.  There are at most \(n\) main node
visits, where \(n\) is the input's node count.  Each of at most \(n\)
preliminary renamings visits at most \(n\) nodes.  Entering or returning
from a frame and moving to the next child require a fixed number of stack
operations, as all constructors have at most two children.  For example
\(20(n+1)^2\) iterations suffice if a frame step performs its tag test,
list projections, variable-list calculation, and bounded fresh search in
one update; each of these updates is already primitive recursive.
After termination make the update leave the final state fixed.  Iterating
for this bound gives the result by the bounded-loop principle.  An induction
on the processed tree, with the preliminary renaming clause just verified,
identifies it with the defined capture-avoiding substitution, including a
renaming even when the substituted variable does not occur freely.

For alpha equivalence, traverse a formula with the list of enclosing binder
indices, most recent first.  Replace each bound variable occurrence by the
position of its first matching entry, each free occurrence by its original
index with a distinct free-variable tag, and each binder by its quantifier
tag alone.  More precisely use \(N_{18}(j)\) for a bound occurrence at
zero-based stack position \(j\), \(N_{17}(i)\) for a free \(x_i\),
and \(N_{11}(b)\) or \(N_{12}(b)\) for a binder with matrix record
\(b\); retain all other constructor tags and child positions.
These are comparison records, not formula codes.  Stack operations and searches
through at most \(n\) enclosing binders give a bounded primitive-recursive
traversal.  Two formulas are alpha-equivalent exactly when these records
agree.  To justify both directions, each generating fresh renaming in the
definition preserves the record, as do equivalence and constructor
congruence.  Conversely, for equal records choose distinct variable indices
for corresponding binder positions, all absent from both original formulas.
Rename binders from outside inward to these indices.  Each step is a
permitted generating fresh renaming; nearest-binder positions and free
indices then make the two resulting formulas literally identical.  There
are finitely many binders, so this proves the converse without an appeal
to semantic equivalence.  Reject nonformulas before performing this test.

\paragraph{Exact derivation records.}
A derivation record is \(N_{15}(R)\), where \(R\) is a nonempty list of
rows in postorder.  Row \(j\) has exactly five fields
\([r,a,J,U,l]\): a rule tag, concluding formula code, ordered list of
premise-row indices, list of rule parameters, and assumption marker.
The marker is \(0\) for an unmarked assumption and \(i+1\) for label
\(i\); it must be zero at other rules.  Every premise index is less than
\(j\), every nonlast row occurs exactly once among all premise lists,
and the last row occurs zero times.  These requirements encode a tree,
not shared subderivations: following parent edges strictly increases the
index and must end at the sole row without a parent.  Thus every row is
connected to the last row.  The empty list is not a derivation.

Assign tags \(0,\ldots,19\), respectively, to
\[
\begin{gathered}
 \mathsf{assumption},\ \land I,\ \land E_1,\ \land E_2,\
 \lor I_1,\ \lor I_2,\ \lor E,\ \to I,\ \to E,\
 \neg I,\ \neg E,\ \bottom E,\ \mathsf{RAA},\\
 \forall E,\ \exists I,\ \forall I,\ \exists E,\ \alpha,\ =I,\ =E.
\end{gathered}
\]
All other rule tags are rejected.  Parameters are empty unless specified
below.  Reject any wrong list length, wrong type, extra parameter, or
failure of a displayed condition.  Every row formula must be a formula.

Maintain at each row the list \(O_j\) of its open assumption-leaf indices.
An assumption has no premises, empty parameters, and \(O_j=[j]\).
For any non-discharge rule take the concatenation of the premise lists;
\(=I\) has no premises and hence the empty open list.  For a list \(O\),
write \(D_i^\theta(O)\) for removal of all leaves whose marker is
\(i+1\) and whose formula is literally \(\theta\).  Whenever a rule
calls this removal, require at least one such leaf.  This implements the
book's nonvacuous, formula-specific labelled discharge; it does not remove
an occurrence in a different premise branch merely because its marker agrees.

The following complete list specifies the row checks.  Letters for formulas
and terms below denote decoded objects, so all equalities in these checks
are literal syntactic equality except the stated alpha test.
\begin{enumerate}
\item \(\land I\) has premises \(\varphi,\psi\) and conclusion
\(\varphi\land\psi\).  The two eliminations select the indicated child
of the sole conjunctive premise.  The two disjunction introductions have
one premise equal to the indicated child of their disjunctive conclusion.
\item \(\lor E\) has parameters \([i,j]\), premises
\(\varphi\lor\psi,\chi,\chi\), and conclusion \(\chi\).  Its open
list is the major premise's open list, followed by
\(D_i^\varphi(O_{\rm second})\) and \(D_j^\psi(O_{\rm third})\).
\item \(\to I\) has parameter \([i]\), premise \(\psi\), and
conclusion \(\varphi\to\psi\); its open list is
\(D_i^\varphi(O_{\rm premise})\).  The \(\to E\) premise order is
\(\varphi,\varphi\to\psi\), with conclusion \(\psi\), as in the
body rule display.  The rule sheet displays the opposite premise order;
exchanging those two child positions is the exact encoding translation.
\item \(\neg I\) has parameter \([i]\), premise \(\bottom\),
conclusion \(\neg\varphi\), and removes \(\varphi^i\).
\(\neg E\) has premises \(\varphi,\neg\varphi\) and concludes
\(\bottom\).  The rule \(\bottom E\) has sole premise \(\bottom\)
and any formula as conclusion.  The rule \(\mathsf{RAA}\) has parameter
\([i]\), premise \(\bottom\), conclusion \(\varphi\), and removes
\((\neg\varphi)^i\).
\item \(\forall E\) has parameters \([i,b,t]\), where \(b\) codes
\(\varphi\) and \(t\) a term; its premise is \(\forall x_i\varphi\)
and conclusion is the computed \(\varphi[t/x_i]\).  The rule
\(\exists I\) uses the same parameters, premise \(\varphi[t/x_i]\),
and conclusion \(\exists x_i\varphi\).
\item \(\forall I\) has parameter \([i]\), premise \(\varphi\),
conclusion \(\forall x_i\varphi\), and requires \(i\) absent from
the free-variable list of every open premise assumption.
\item \(\exists E\) has parameters \([i,b,y,k]\), where \(b\)
codes \(\varphi\); its premises are \(\exists x_i\varphi\) and
\(\psi\), and its conclusion is \(\psi\).  Remove
\(D_k^{\varphi[x_y/x_i]}\) from the minor open list and concatenate
the remaining list with the major open list.  Require \(y\) absent
from the all-variable lists of \(\psi\), \(\exists x_i\varphi\),
and every remaining open assumption.  This is the source's absence
condition, which is stronger than merely not being free.
\item \(\alpha\) has one premise alpha-equivalent to its conclusion.
The rule \(=I\) has parameter \([t]\), no premises, and conclusion
\(t=t\).  The rule \(=E\) has parameters \([r,t,b,i]\), where
\(b\) codes \(\varphi\), premises \(r=t,\varphi[r/x_i]\), and
conclusion \(\varphi[t/x_i]\).  No occurrence of \(x_i\) in
\(\varphi\) is required.
\end{enumerate}
There is no additional repetition, arbitrary equivalence, logical-axiom,
or implicit substitution rule.  A use of a universally quantified
nonlogical axiom is an assumption of that sentence followed by the actual
\(\forall E\) steps.  Thus for a fixed finite axiom set
\(T=\{\tau_0,\ldots,\tau_{m-1}\}\), the axiom-leaf test is the
finite disjunction \(a=\ulcorner\tau_j\urcorner\), \(j<m\).
Its instances with supplied terms are checked by the substitution and
\(\forall E\) tests just specified, not treated as a new axiom schema.
This also covers \(m=0\): no open assumption is then permitted.  The
universal closures in the definition of \(\mathsf Q\), for example,
are its seven axiom sentences, rather than an unspecified collection of
logical or substitution axioms.

\begin{definition}[Computably enumerable theory]
Relative to the fixed G\"odel numbering, a theory
\(T\subseteq\Sent(\Lang_A)\) is \emph{computably enumerable} iff
\[
  \set{\ulcorner\sigma\urcorner:\sigma\in T}
\]
is a computably enumerable set of natural numbers.  This term classifies the
given set of axioms \(T\); it does not mean that \(T\) is deductively closed,
and it does not define \(T\) as the set of its derivable consequences.  The
distinction is required because the proof predicate below verifies uses of
members of \(T\) by finite accepting-computation certificates.  Fix a
program index \(e_T\) in the existing Part IV numbering whose partial
function has domain exactly this set of axiom codes.  This is the
domain convention for a c.e. set; an enumeration presentation gives such
an index through \cref{thm:enumeration}.
\end{definition}

\paragraph{Certified axiom leaves.}
Write \(V(e,a,s)\) for the numerical truth value of the already fixed
bounded formula \(\mathsf{Comp}(\overline e,\overline a,\overline s)\)
in \(\mathcal N\).  The arguments here are numbers; \(V\) is a
metamathematical characteristic function, not an added arithmetic symbol.
It is primitive recursive.  Indeed evaluate each fixed term by its
\(0,S,+,\times\) operations, an atomic equality by numerical comparison,
and each Boolean connective by its Boolean truth table.  For a marked
bounded existential use a bounded sum followed by a positivity test, and
for a marked bounded universal use a bounded product.  Substitute the
primitive-recursively evaluated bound into that recursion.  Its bound
variable is absent from the bounding term.  Empty existential ranges give
zero, and empty universal ranges give one.  Structural induction on this
one fixed formula in its accepted marked-bounded presentation, expanding
its relational abbreviations while retaining the marked bounds for this
evaluation, proves the claim.  The meanings of the bounded-quantifier
abbreviations identify this value with truth of their full first-order
expansion.  This argument asserts no uniform truth predicate for
arbitrary unbounded arithmetic formulas.

A \(T\)-proof code is \(N_{16}(d,W)\), where \(d=N_{15}(R)\)
passes all the derivation checks and has a sentence as root formula.
If its root open list is \([j_0,\ldots,j_{u-1}]\), require
\(W=[s_0,\ldots,s_{u-1}]\) of exactly that length.  For each \(v<u\),
the leaf formula \(a_{j_v}\) must be a sentence and satisfy
\(V(e_T,a_{j_v},s_v)=1\).  Discharged temporary assumptions need no
axiom certificate.  When \(u=0\), \(W=0\) is required and there is
no call to an axiom-membership test.  Define the numerical relation
\(\mathsf{Proof}_T(p,x)\) by these checks together with equality of
\(x\) and the root formula code.  Every other number is rejected.
For a fixed finite \(T\), one may instead use \(N_{16}(d,0)\) and
apply its finite membership test to all root open leaves; the selected
convention must remain fixed when naming a proof predicate.

\begin{example}[A finite-sequence code]
One illustrative component code assigns to
\((a_0,\ldots,a_{k-1})\) the natural number
\[
  \operatorname{pair}_{\N}\left(k,\prod_{i<k}p_i^{a_i+1}\right),
\]
where \(p_i\) is the \(i\)-th prime.  For \((2,0,1)\), the product is
\[
  2^3\cdot3^1\cdot5^2=600,
\]
so the complete code is \(\operatorname{pair}_{\N}(3,600)\).  The
exponent \(a_i+1\) rather than \(a_i\) is what distinguishes trailing
zeros; pairing with \(k\) additionally makes the length available by a
projection rather than by factoring.  This is a separate prime-power
illustration; the governing syntax and derivation codes in this section
are the constructor/list codes above, and the computation witnesses retain
the Part IV \(\beta\) coding.
\end{example}

\begin{theorem}[Arithmetization of syntax]
For the explicit numbering above, symbol-kind, term, formula, and sentence
recognition; conversions between canonical strings and trees; numeral
formation; capture-avoiding substitution (in particular numeral
substitution); alpha equivalence; finite-axiom membership and explicitly
specified axiom instantiations; and correctness of the specified finite
\(\ND_{\FOL}\) derivation records are primitive recursive.  Correctness
from a fixed finite set of nonlogical axioms is primitive recursive.
For every fixed c.e. axiom set \(T\), the certified relation
\(\mathsf{Proof}_T\) is primitive recursive and, for every sentence
\(\sigma\),
\[
 T\derives\sigma
 \quad\Longleftrightarrow\quad
 \exists p\in\N\;\mathsf{Proof}_T(p,\ulcorner\sigma\urcorner).
\]
On the right is the numerical relation just defined.  This is a
standard-model coding theorem, without a claim about what \(\mathsf Q\)
proves concerning all codes.
\end{theorem}

\begin{proof}
The list inverses, child-code recursion, string bound, substitution
traversals, and alpha-record argument above prove the assertions about
syntax, including their total defaults.  In particular none follows merely
from termination of an unspecified search.

For a derivation code, loop through its \(\ell(R)\) rows.  Each row has
fixed length, each premise list is checked for the required rule arity,
and each reference is earlier than its row.  Parent counts are bounded
sums over these lists.  Compute \(O_j\) from already computed rows,
using list concatenation and the indicated discharge filters.  Each local
rule test is a finite combination of constructor decoding, the proved
substitution and alpha tests, and a loop over the current open list for
freshness or discharge.  These are all primitive recursive.  Store the
open lists in the accumulated row table, so the whole computation is
primitive recursion on \(\ell(R)\), followed by its root check.
Malformed or empty records return false.  Finite-axiom membership adds
only its finite disjunction; certified c.e. membership adds a loop of
length \(\ell(W)\) through the primitive-recursive function
\(V(e_T,a,s)\).  Thus both versions of the proof relation are primitive
recursive, even though membership in a c.e. axiom set need not be decidable.

It remains to prove that the checker recognizes the intended derivations.
Induct over the postorder rows.  An assumption row is exactly an open
assumption leaf; an \(=I\) row is exactly the premise-free identity rule.
At each other row the child subtrees are disjoint, the checked formula
pattern is the corresponding displayed rule, and concatenation lists
exactly their open assumption occurrences.  Each discharge filter removes
exactly the marked occurrences specified for its premise and formula and
checks nonvacuity.  The remaining lists therefore are precisely the open
assumptions of the new subtree.  The \(\forall I\) free-variable test
and the \(\exists E\) all-occurrence test are exactly their respective
side conditions.  The substitution and alpha tests have already been
identified with the book's syntactic operations.  This proves the
induction invariant, not just matching of the concluding formulas.
Connectedness makes the last row's subtree the whole recorded tree.

Suppose now that \(\mathsf{Proof}_T(p,\ulcorner\sigma\urcorner)\)
holds.  The preceding induction extracts a genuine
\(\ND_{\FOL}\) derivation of \(\sigma\).  For every open leaf,
\cref{thm:bounded-verification} turns its successful \(V\) test into
an accepting numerical computation of program \(e_T\) on its formula
code.  By the choice of \(e_T\) that sentence belongs to \(T\).
Consequently the extracted derivation proves \(T\derives\sigma\).

Conversely, take a finite \(\ND_{\FOL}\) derivation of \(\sigma\)
with open assumptions in \(T\).  List its actual tree nodes in postorder,
put the displayed rule witnesses (terms, matrices, indices, discharge
labels) into their parameter lists, and use the existing assumption labels
with the stated marker convention.  The rule-by-rule induction in the
other direction shows that all row and open-list checks succeed.  There
are finitely many root open leaves.  For each, its membership in \(T\)
means that \(\varphi_{e_T}\) is defined on its code; the other direction
of \cref{thm:bounded-verification} provides a finite witness \(s_v\).
Choosing witnesses for this finite list requires only finite successive
existential instantiation.  Their list \(W\), or the empty list when
there are no open leaves, produces the required proof code.  The finite
axiom version has the identical argument with direct membership checks.
\end{proof}

\begin{definition}[Proof and provability predicates for the certified presentation]
\label{def:fixed-proof-predicate}
Let \(T\) be computably enumerable and contain \(\mathsf Q\).  Use the
certified primitive-recursive numerical relation \(\mathsf{Proof}_T(p,x)\)
just constructed, with its fixed index \(e_T\).  By the
\(\Sigma^0_1\) clause of \cref{thm:representability}, fix the same symbol
for a
specific \(\Sigma^0_1\) formula that numeralwise represents this relation in
\(\mathsf Q\).  Define
\[
  \mathsf{Prov}_T(x)\defeq\exists p\,\mathsf{Proof}_T(p,x).
\]
The declared \(\Sigma^0_1\) class is closed under unbounded existential
quantification, so \(\mathsf{Prov}_T\) is a \(\Sigma^0_1\) formula by
the syntactic hierarchy definition.  The
\(\beta\) coding of \cref{lem:goedel-beta} supplies the computation
histories inside the axiom certificates; the surrounding syntax trees
and proof sequences use the explicitly different list codes above.
The numeralwise representing formula is supplied by
\cref{thm:representability}, in addition to the numerical checker theorem.
This definition retains the certified presentation and its chosen formula.
The standard \(\PA\) development separately specifies the direct formula
\(\mathsf{Prv}^{19}_{\PA}\) in
\cref{def:p05-selected-direct-pa-graph}; its internal derivability
conditions are proved in
\cref{thm:p05-direct-pa-internal-mp,thm:p05-direct-pa-d3}.
Those constructions do not identify an arbitrary numeralwise representative
with the direct formula inside arithmetic.
\end{definition}

\begin{lemma}[The first derivability condition from numeralwise representation]
\label{lem:p05-numeralwise-d1}
For every computably enumerable \(T\supseteq\mathsf Q\) and every
sentence \(\sigma\), the predicate just defined satisfies
\[
 T\derives\sigma\quad\Longrightarrow\quad
 T\derives\mathsf{Prov}_T(\overline{\ulcorner\sigma\urcorner}).
\]
\end{lemma}
\begin{proof}
Suppose \(T\derives\sigma\).  The numerical direction of
\cref{thm:arithmetization-syntax} supplies a standard number \(p\) for
which the numerical relation
\(\mathsf{Proof}_T(p,\ulcorner\sigma\urcorner)\) holds.
The positive clause of numeralwise representation therefore gives an actual
\(\mathsf Q\)-derivation of the arithmetic formula
\[
 \mathsf{Proof}_T(\bar p,\overline{\ulcorner\sigma\urcorner}).
\]
Every open axiom of that derivation belongs to \(T\), since
\(\mathsf Q\subseteq T\).  It is consequently a \(T\)-derivation.
Apply the actual \(\exists I\) rule with the closed term \(\bar p\)
to obtain the required \(\mathsf{Prov}_T\) instance.  This argument uses
the cited numeralwise representation clause in addition to the numerical
checker theorem; it proves the displayed metatheoretic implication.
\end{proof}

\begin{warning}[Mentioning a sentence inside arithmetic]
The expression \(\mathsf{Prov}_T(\varphi)\) is ill-typed if \(\varphi\) is a
formula rather than a number term.  The correct expression is
\(\mathsf{Prov}_T(\overline{\ulcorner\varphi\urcorner})\), often abbreviated
only after the coding convention has been declared.
\end{warning}

\begin{exercise}
\label{ex:arithmetic-coding-02}
Define a primitive recursive relation saying that a finite sequence is a
correct addition computation
\(a,a+1,\ldots,a+b\).  Explain how the relation supports representability of
addition.
\end{exercise}

\begin{exercise}
\label{ex:arithmetic-coding-03}
Distinguish the following three assertions by giving their types:
\(T\derives\sigma\),
\(\mathcal N\entails\mathsf{Prov}_T(\overline{\ulcorner\sigma\urcorner})\),
and
\(T\derives\mathsf{Prov}_T(\overline{\ulcorner\sigma\urcorner})\).
State one hypothesis that connects each adjacent pair.
\end{exercise}
